\clearpage
\section{Implementation}

This chapter maps the design decisions to the final Java implementation. It
focuses on the main mechanisms that realize ownership, synchronization, and
state publication.

\subsection{Runtime and View Integration}

The code is divided by responsibility:

\begin{center}
\begin{tabular}{|>{\bfseries}p{3.2cm}|p{10.2cm}|}
    \hline
    Package or class & Implementation responsibility \\
    \hline
    Playable launchers
        & \texttt{SequentialPoool}, \texttt{ThreadedPoool}, and \texttt{TaskBasedPoool} assemble the selected model, runtime, and view; copy snapshots; handle restart and rendering. \\
    \hline
    \texttt{runtime}
        & Serialize commands, execute controller ticks, publish snapshots, and run the bot. \\
    \hline
    \texttt{model.game}
        & Enforce turns, shooting rules, score updates, and match termination. \\
    \hline
    \texttt{model.physics}
        & Store the board and implement sequential, platform-thread, and task-based physics. \\
    \hline
    \texttt{view}
        & Convert input into callbacks and render synchronized copies of the game state. \\
    \hline
\end{tabular}
\end{center}

The sequential launcher owns the model directly and advances it in the same
loop that updates the view. The concurrent launchers instead own a
\texttt{GameRuntime}. Their rendering loops only submit commands and consume
published state.

Both concurrent runners delegate game semantics to \texttt{GameLoop}. The
runner chooses where a tick executes, but the tick protocol is shared:

\begin{lstlisting}[style=codestyle,caption={Shared single-writer controller tick, from \texttt{runtime.GameLoop}},label={lst:shared-game-loop}]
public void tick(long tickMillis) {
    drainCommands();
    if (!game.snapshot().isFinished()) {
        game.step(tickMillis);
    }
    publishSnapshot();
}
\end{lstlisting}

\texttt{ThreadedGameRunner} calls this method from a dedicated platform
thread. \texttt{TaskBasedGameRunner} schedules it at a fixed rate on a
single-thread \texttt{ScheduledExecutorService}. In both cases,
\textbf{only that controller execution advances \texttt{GameModel}}.

\texttt{CommandMailbox} is the monitor between asynchronous producers and the
model owner. Swing and \texttt{BotAgent} submit a
\texttt{Function<GameModel,T>} and receive a \texttt{CompletableFuture<T>}.
Submission only appends to a FIFO queue; execution occurs when
\texttt{GameLoop} drains that queue at the beginning of the next tick.

After commands and physics have completed, the loop creates the next published
state. The resulting \texttt{RuntimeGameSnapshot} contains the immutable logical
\texttt{GameSnapshot}, copied ball records, holes, and the bot preview vector.

\begin{center}
\textbf{FIFO commands} $\rightarrow$ \textbf{single model writer}
$\rightarrow$ \textbf{copied runtime snapshot}
\end{center}

The concurrent launchers repeatedly read the latest
\texttt{RuntimeGameSnapshot} and copy it into \texttt{ViewModel}. All
\texttt{ViewModel} updates and getters are synchronized, and its public
rendering records contain only positions, radii, and preview information.
\texttt{ViewFrame} therefore never traverses the live \texttt{Board}.

Mouse and keyboard listeners remain inside \texttt{ViewFrame}. They calculate
the shot direction and intensity, then invoke the callback supplied by the
launcher. In a concurrent version that callback calls
\texttt{GameRuntime.shootHuman(...)}; the runtime places the request in
\texttt{CommandMailbox}. Input validation and match progression consequently
remain on the controller owner.

Rendering remains separated from simulation: Swing consumes the copied
\texttt{ViewModel} state and never runs physics on the EDT.

\subsection{Physics Step Structure}

All physics engines implement \texttt{PhysicsStepper}. A concurrent engine
holds the board monitor for the entire step and divides long intervals into
bounded sub-steps:

\begin{lstlisting}[style=codestyle,caption={Serialized board ownership and bounded sub-steps, from \texttt{model.physics.threaded.ThreadedPhysicsEngine}},label={lst:physics-step-entry}]
synchronized (board) {
    long remaining = elapsedMillis;
    while (remaining > 0) {
        long dt = Math.min(maxStepMillis, remaining);
        stepOnce(board, dt);
        remaining -= dt;
    }
}
\end{lstlisting}

The board lock makes each sub-step atomic to external observers. Parallelism is
internal to \texttt{stepOnce}: workers may update disjoint portions of the
board state, but only while that physics step owns the board.

Independent collections are split into contiguous ranges. Unless a fast path
selects one range, $R=\min(W,N)$ workers or tasks process $N$ items. Each range
receives $\lfloor N/R \rfloor$ items and the first $N \bmod R$ ranges receive
one additional item. This guarantees:

\begin{itemize}
    \item no item is processed twice or omitted;
    \item range sizes differ by at most one;
    \item each worker writes only to its assigned slice or private result.
\end{itemize}

The same calculation partitions active balls, occupied cells, and touched
ball indexes.
\subsection{Threaded and Task-Based Execution}

\texttt{ThreadedPhysicsEngine} owns a fixed array of long-lived
\texttt{PhysicsWorker} objects. Parallel phases are distributed to these
workers and synchronized through explicit completion barriers.

\begin{lstlisting}[style=codestyle,caption={Platform-thread dispatch and barrier, adapted from \texttt{model.physics.threaded.ThreadedPhysicsEngine}},label={lst:threaded-dispatch}]
var completion = new WorkerCompletionMonitor(activeWorkerCount);
workers[workerIndex].assign(() -> {
    for (int i = rangeStart; i < rangeEnd; i++) {
        activeBalls.get(i).updateState(dt, bounds);
    }
}, completion);
completion.await();
\end{lstlisting}

\texttt{TaskBasedPhysicsEngine} owns a fixed \texttt{ExecutorService}. A
parallel phase submits one independent \texttt{Runnable} per range, and
\texttt{Future.get()} is used as the phase barrier.

\begin{lstlisting}[style=codestyle,caption={Executor submission and Future barrier, adapted from \texttt{model.physics.taskbased.TaskBasedPhysicsEngine}},label={lst:task-dispatch}]
futures.add(executor.submit(() -> {
    for (int i = rangeStart; i < rangeEnd; i++) {
        activeBalls.get(i).updateState(dt, bounds);
    }
}));
awaitAll(futures);
\end{lstlisting}

\subsection{Collision Pipeline}

Each fixed sub-step uses the same structure in both concurrent engines:
parallel movement on disjoint ranges, sequential hole handling, parallel broad
and narrow collision detection on private or disjoint data, deterministic merge
of collision contributions, and final application of one accumulated correction
per touched ball.

The uniform grid acts as the broad phase. Each worker builds private data,
collects collision pairs, and leaves the coordinator responsible for the
ordered merge and accumulation. The final apply phase may update the actual
board state in parallel, but only over disjoint touched-ball ranges. This
avoids conflicting writes while keeping the physics result aligned with the
sequential semantics.

\begin{center}
\textbf{private detection, ordered accumulation, disjoint application}
\end{center}

Shutdown follows resource ownership. A runner stops ticks and closes the
mailbox. The threaded version joins controller and bot threads, then closes its
physics workers. The task-based version terminates its executors and finally
shuts down the physics pool.

Benchmark classes are development-only and are excluded from the delivery
package. The submitted archive contains the Maven descriptor, production
sources, tests, and this report PDF.
